% #############################################################################
% This is the Data chapter
% !TEX root = ../main.tex
% #############################################################################
\fancychapter{Data}
\label{chap:data}
\label{sec:data}

The reported experiments use four main resources with different roles. Each
resource is first converted into a common aligned-pair format containing a BP
sentence, an EP sentence, optional dataset-specific metadata, and a split label.
This chapter describes the origin and preparation of each resource.
\Cref{tab:data-resources} summarizes aligned EP/BP pair counts for each resource
before task-specific rendering.

\begin{table}[htbp]
  \centering
  \thesistablefont
  \caption{Aligned EP/BP pair counts by resource and split before task-specific rendering.}
  \label{tab:data-resources}
\begin{tabular}{@{}lccc@{}}
    \toprule
    Resource & Train & Validation & Test \\
    \midrule
    OpenSubtitles & 10,347,883 & 0 & 0 \\
    GPT-Wikipedia & 4,519 & 564 & 566 \\
    FRMT & 2,502 & 20 & 2,611 \\
    Golden Collection & 0 & 0 & 500 \\
    \bottomrule
  \end{tabular}

\end{table}

\section{OpenSubtitles}
\label{sec:data-opensubtitles}

OPUS OpenSubtitles \citep{4992de1b5fb34f3e9691772606b36edf} is a large
collection of aligned subtitles from movies and television series. It provides
the scale needed for an initial broad adaptation phase and exposes the models
to diverse topics and informal dialogue. However, its scale comes with
limitations, as subtitle segments can be elliptical, weakly aligned, or
stylistically distant from conservative EP/BP adaptation. Its role in the
experiments is broad exposure rather than precise control of
individual variety markers.

{\tolerance=1500\looseness=-1
After exact duplicate pairs are removed, preprocessing first removes
subtitle-specific formatting noise, including leading dialogue dashes, speaker
labels, dash artifacts after punctuation, irregular whitespace, wrapping and
unbalanced quotation marks, and fully capitalized line formatting. The next step
repairs segmentation when adjacent subtitle rows appear to belong to the same
sentence. Rows \(i\) and \(i+1\) are merged when either the EP or BP side
provides a continuation cue: row \(i\) ends in a comma, semicolon, colon, or dash,
or contains fewer than 40 characters, and the first alphabetic character of row
\(i+1\) is lowercase. A separate merge condition handles cases where the paired
sentences show inconsistent apparent sentence boundaries. Whenever a merge
condition is met on either side, both BP and EP rows are joined so their
alignment is preserved. \Cref{tab:opensubs-merge-examples} illustrates
representative segmentation repairs. In each example, ``Line \(i\)'' and ``Line
\(i+1\)'' are adjacent subtitle rows, while ``Merged sentence'' shows the resulting BP text.\par}

\begin{table}[htbp]
  \centering
  \thesistablefont
  \caption{Examples of OpenSubtitles segmentation repairs using the BP side of each aligned EP/BP pair.}
  \label{tab:opensubs-merge-examples}
  \begin{tabularx}{\textwidth}{@{}YYYY@{}}
    \toprule
    Merge cue & Line \(i\) & Line \(i+1\) & Merged sentence \\
    \midrule
    Unfinished punctuation + lowercase continuation &
    \textit{N\~ao achava que o basquete era assim, o cansa\c{c}o, a vit\'oria,} &
    \textit{o cansa\c{c}o.} &
    \textit{N\~ao achava que o basquete era assim, o cansa\c{c}o, a vit\'oria, o cansa\c{c}o.} \\
    \addlinespace
    Short fragment split across rows &
    \textit{Massa Macinaia.} &
    \textit{Deve ser uma vila perto de...} &
    \textit{Massa Macinaia. Deve ser uma vila perto de...} \\
    \addlinespace
    Ellipsis + lowercase continuation &
    \textit{N\~ao pode colocar na sua cabe\c{c}a? ...} &
    \textit{eu n\~ao morri na primeira vez.} &
    \textit{N\~ao pode colocar na sua cabe\c{c}a?... Eu n\~ao morri na primeira vez.} \\
    \bottomrule
  \end{tabularx}
\end{table}

{\tolerance=1500\looseness=-1
Surface overlap alone is not a reliable filter for aligned EP/BP pairs, since
valid correspondences may preserve the same content while using different
lexical forms or one-to-many constructions. For example, a BP gerund can
correspond to an EP infinitival construction. The pairs are therefore filtered
with SimAlign \citep{jalili-sabet-etal-2020-simalign}, which uses contextual
token embeddings to provide a bilingual alignment signal rather than relying
on exact word matching.\par}

{\tolerance=1500\looseness=-1
Using SimAlign requires choosing how the alignments are obtained. The first
choice is the multilingual encoder used to compute contextual embeddings for
both sentences. The second is the granularity of the alignment, since token
connections can be produced over subword pieces or over the original words
after combining the representations of their subword pieces.
The third is the alignment extraction method that converts token similarities
into links between the sentences of the pair. In this work, the filter uses the
XLM-RoBERTa (XLM-R)-based encoder configuration, word-level alignments, and the
\texttt{inter} extraction method. This extraction can be understood as a
similarity matrix whose rows are BP tokens and whose columns are EP tokens.
Each cell estimates how close two contextual token representations are, and
\texttt{inter} keeps a link only when the two tokens select each other as their
highest-scoring option. The resulting links are used only as evidence that the
paired sentences still express the same content. They are converted into a coverage
score over normalized word tokens. Shorter subtitle fragments are allowed to
have lower coverage, while longer pairs must show stronger bilingual
alignment. A separate length-ratio check rejects pairs where one sentence is
much shorter or longer than its counterpart. The exact scoring rule, thresholds,
alternative SimAlign options, and an illustrative token-similarity matrix are detailed in
\Cref{chap:appendix-opensubs-filtering}.\par}

\section{GPT-Wikipedia}
\label{sec:data-gpt-wikipedia}

{\tolerance=1500\looseness=-1
GPT-Wikipedia is the main constructed resource for controlled EP/BP
rewriting. It was created to provide more conservative supervision than could
be expected from subtitle data alone, while retaining broader topical coverage
than a small manually curated benchmark. The resource is designed around
near-literal EP/BP pairs, where the two sentences remain close in meaning and
word overlap. Stricter
conservative subsets are also derived from this resource for later
reward-based continuation experiments.\par}

GPT-Wikipedia is built from Portuguese Wikipedia article dumps used as
inspiration material. For each article, only the first few sentences from
informative paragraphs are kept. Paragraphs that are too short, list-like, or
punctuation-poor are discarded. These texts are used only to provide topics and
contexts for generation, not as aligned examples.

Generated pairs are obtained with a Portuguese prompt sent to OpenAI GPT-4o
through the IAedu API. The
template requests new EP/BP pairs in JSON format and favors minimal lexical
change. Its main constraints are summarized by the translated and simplified
excerpt below.

\begin{quote}
\small
Generate \(N\) new bilingual examples in valid JSON, with one
EP sentence and one BP sentence per example. Use the
reference passages only as inspiration for topic, domain, and register. Keep
the two variants as close as possible. Do not write free paraphrases or change
meaning without need. Change only clear lexical or morphosyntactic contrasts.
In non-identical pairs, prefer one or two clear differences and avoid extensive
rewrites.
\end{quote}

The generated batches are then merged and checked for suspicious differences.
The check compares observed edits with an allowlist of expected EP/BP
correspondences and flags unexpected lexical substitutions or broader reordering
for review.

{\tolerance=1500\looseness=-1
During review, variety-relevant differences are kept with the EP expression in
the EP sentence and the BP expression in the BP sentence. Differences judged unnecessary
for EP/BP adaptation are collapsed to the same wording. As a result, the final
cleaned resource is not restricted to non-equal pairs. The dataset contains
1,648 equal pairs among 5,649 pairs after cleanup.\par}

\section{FRMT}
\label{sec:data-frmt}

{\tolerance=1500\looseness=-1
FRMT \citep{riley2023frmtbenchmarkfewshotregionaware} is a benchmark for
few-shot region-aware MT. It provides English sentences translated into
regional varieties, including EP and BP. Its examples are organized into three
buckets. The lexical bucket is designed around
articles likely to elicit region-dependent terminology. The entity
bucket focuses on people, places, institutions, and other region-associated
entities. The random bucket draws from randomly selected high-quality
Wikipedia articles and evaluates broader regional behavior without
preselecting a lexical or entity contrast.\par}

FRMT is not generated internally, but it is converted into the same pair
format used by the rest of the resources so that training and evaluation
scripts can share a common interface. Its original bucket labels are retained
as metadata, although the reported evaluation aggregates over the complete
FRMT split.

{\tolerance=1500\looseness=-1
The complete FRMT training split provides additional supervised data in the
GPT-Wikipedia + FRMT condition, while the held-out test split is used for evaluation.
A stricter conservative subset is selected later for reward-based continuation,
as described in Section~\ref{sec:method-stage-c}.\par}

\section{Golden Collection}
\label{sec:data-golden-collection}

The Golden Collection is a manually curated evaluation resource introduced by
\citet{sanches2024brazilianportugueseeuropeanportuguese}. It contains 500 BP
sentences divided across Science, Legislative, Literature, Social Media, and
Wiki topics. In the original dataset, each source is associated with EP references, including a
manually curated reference and an automatic DeepL reference.

{\tolerance=1500\looseness=-1
Each BP sentence is paired with the manually curated EP
reference, rather than the automatic DeepL reference. The resulting 500 EP/BP
pairs are used only for evaluation. For rewriting evaluation, each pair is
converted into two examples: one BP-to-EP and one EP-to-BP. When both
sentences are identical, the expected output is the unchanged sentence.
\par}
